\documentclass[10pt,a4paper]{article}
\usepackage[a4paper,margin=14mm,top=12mm,bottom=14mm]{geometry}
\usepackage{fontspec}
\setmainfont{Liberation Sans}
\usepackage{polyglossia}
\setdefaultlanguage{french}
\usepackage{amsmath}
\usepackage{xcolor}
\usepackage{graphicx}
\usepackage[most]{tcolorbox}
\usepackage{tabularx,booktabs,array,multirow}
\usepackage{enumitem}
\usepackage{fancyhdr}
\usepackage{lastpage}
\usepackage[hidelinks]{hyperref}
\definecolor{uimm}{HTML}{005677}
\definecolor{cyan}{HTML}{13A2DD}
\definecolor{fond}{HTML}{E6F2FA}
\definecolor{ok}{HTML}{2E7D32}
\definecolor{ko}{HTML}{B3261E}
\definecolor{violet}{HTML}{7B1FA2}
\setlength{\parindent}{0pt}
\setlength{\parskip}{3pt}
\setlist{nosep}
\setlist[itemize]{leftmargin=13pt,label=\textcolor{cyan}{\textbullet}}
\newcommand{\titre}[1]{\vspace{5pt}{\color{uimm}\large\bfseries #1}\par\vspace{-3pt}{\color{cyan}\rule{\linewidth}{1.2pt}}\par\vspace{0pt}}
\tcbset{boite/.style={colback=fond,colframe=uimm,boxrule=0.6pt,arc=2pt,left=6pt,right=6pt,top=4pt,bottom=4pt},
        defi/.style={colback=violet!6,colframe=violet,boxrule=0.8pt,arc=2pt,left=6pt,right=6pt,top=4pt,bottom=4pt,
                     title={\bfseries Défi},coltitle=white,colbacktitle=violet,fonttitle=\small}}
\newcommand{\bandeau}[3]{%
\begin{tcolorbox}[colback=uimm,colframe=uimm,arc=2pt,left=4pt,right=8pt,top=4pt,bottom=4pt]
\begin{minipage}[c]{18mm}\setlength{\fboxsep}{1.2pt}\fcolorbox{white}{white}{\includegraphics[height=16.5mm]{logo_uimm.jpg}}\end{minipage}\hspace{4mm}%
\begin{minipage}[c]{\dimexpr\linewidth-22mm\relax}\color{white}
{\small #1}\\[1pt]{\LARGE\bfseries #2}\\[2pt]{\large #3}\hfill{\small CQPM Technicien qualité}
\end{minipage}
\end{tcolorbox}}
\newcommand{\reponse}[1][1]{\par\vspace{2pt}{\color{gray!60}\foreach \i in {1,...,#1}{\rule{\linewidth}{0.3pt}\par\vspace{9pt}}}}
\usepackage{pgffor}
\pagestyle{fancy}\fancyhf{}
\renewcommand{\headrulewidth}{0pt}
\newcommand{\pied}[1]{\fancyfoot[L]{\footnotesize\color{uimm}#1}\fancyfoot[R]{\footnotesize\color{uimm}Page \thepage{} sur \pageref{LastPage}}}

\usepackage{tikz}
\usepackage{longtable}
\renewcommand{\bandeau}[3]{%
\begin{tcolorbox}[colback=uimm,colframe=uimm,arc=2pt,left=4pt,right=8pt,top=4pt,bottom=4pt]
\begin{minipage}[c]{18mm}\setlength{\fboxsep}{1.2pt}\fcolorbox{white}{white}{\includegraphics[height=16.5mm]{logo_uimm.jpg}}\end{minipage}\hspace{4mm}%
\begin{minipage}[c]{\dimexpr\linewidth-22mm\relax}\color{white}
{\small #1}\\[1pt]{\LARGE\bfseries #2}\\[2pt]{\large #3}
\end{minipage}
\end{tcolorbox}}
% en-tête d'une fiche : code, titre, question, bloc, durée, organisation, niveau
\newcommand{\entete}[7]{%
\begin{tcolorbox}[colback=uimm,colframe=uimm,arc=2pt,left=4pt,right=6pt,top=4pt,bottom=4pt]
\begin{minipage}[c]{18mm}\setlength{\fboxsep}{1.2pt}\fcolorbox{white}{white}{\includegraphics[height=16.5mm]{logo_uimm.jpg}}\end{minipage}\hspace{3mm}%
\begin{minipage}[c]{\dimexpr\linewidth-58mm\relax}\color{white}\raggedright
{\small Base SPC · #4}\\[1pt]{\Large\bfseries #1 · #2}\\[2pt]{\itshape #3}
\end{minipage}\hfill
\begin{minipage}[c]{33mm}\color{white}\small\raggedleft
Durée : \textbf{#5}\\ \textbf{#6}\\ Niveau #7 sur 3
\end{minipage}
\end{tcolorbox}}
\newcommand{\tagE}{{\setlength{\fboxsep}{1.5pt}\colorbox{cyan}{\color{white}\scriptsize\bfseries Ellistat}}}
\newcommand{\tagW}{{\setlength{\fboxsep}{1.5pt}\colorbox{uimm}{\color{white}\scriptsize\bfseries Page web}}}
\tcbset{pred/.style={colback=orange!6,colframe=orange!80!black,boxrule=0.6pt,arc=2pt,left=6pt,right=6pt,top=3pt,bottom=1pt,
                    title={\bfseries Je prédis, avant de commencer},coltitle=white,colbacktitle=orange!80!black,fonttitle=\small},
        retenir/.style={colback=fond,colframe=uimm,boxrule=0.6pt,arc=2pt,left=6pt,right=6pt,top=3pt,bottom=3pt,
                    title={\bfseries À retenir},coltitle=white,colbacktitle=uimm,fonttitle=\small},
        plus/.style={colback=violet!6,colframe=violet,boxrule=0.6pt,arc=2pt,left=6pt,right=6pt,top=3pt,bottom=3pt,
                    title={\bfseries Pour aller plus loin},coltitle=white,colbacktitle=violet,fonttitle=\small},
        act/.style={colback=uimm!10,colframe=uimm,boxrule=0.5pt,arc=2pt,left=5pt,right=5pt,top=2pt,bottom=2pt,before skip=7pt,after skip=3pt},
        verif/.style={colback=orange!6,colframe=orange!80!black,boxrule=0.6pt,arc=2pt,left=6pt,right=6pt,top=3pt,bottom=3pt,
                    title={\bfseries À vérifier sur poste avant la première séance},coltitle=white,colbacktitle=orange!80!black,fonttitle=\small}}
\newcommand{\fichier}[1]{\texttt{\small #1}}
\newcommand{\menu}[1]{\textsf{\textbf{#1}}}
\newcommand{\champ}[1]{\rule{#1}{0.3pt}}

\pied{Base SPC · fiche M8 · Juste partout ?}
\begin{document}
\thispagestyle{fancy}
\entete{M8}{Juste partout ?}{Mon micromètre est-il aussi juste à 5 mm qu'à 25 mm ?}{Le moyen de mesure}{45 min}{en binôme}{3}
\begin{tcolorbox}[boite]\textbf{Le but.} Mesurer cinq cales de longueurs différentes et voir si l'erreur dépend de la longueur.\par\textbf{Il vous faut.} Un micromètre 0-25 mm · Cinq cales : 5, 10, 15, 20 et 25 mm\end{tcolorbox}
\begin{tcolorbox}[pred]Un micromètre juste sur la cale de 10 mm l'est-il forcément sur celle de 25 mm ?\reponse\end{tcolorbox}
\titre{À faire}
\begin{enumerate}[label=\textcolor{cyan}{\textbf{\arabic*.}},leftmargin=15pt,itemsep=1.5pt]
\item Mesurer chaque cale 10 fois, au cliquet, en la retirant à chaque fois. Remplir le tableau.
\item Calculer le biais de chaque cale : moyenne des 10 mesures moins valeur de la cale.
\item \tagE\ \menu{Gage R\&R} > \menu{Linéarité} : 5 produits, 10 valeurs par produit. Pour chaque cale, saisir sa valeur vraie puis ses 10 mesures.
\item \tagE\ Lire la pente, la justesse, et ce qu'Ellistat en dit (significatif ou non).
\end{enumerate}
\vspace{3pt}{\renewcommand{\arraystretch}{1.45}\small\begin{tabularx}{\linewidth}{|l|*{10}{>{\centering\arraybackslash}X|}p{15mm}|p{15mm}|}\hline
\textbf{Cale} & \scriptsize 1 & \scriptsize 2 & \scriptsize 3 & \scriptsize 4 & \scriptsize 5 & \scriptsize 6 & \scriptsize 7 & \scriptsize 8 & \scriptsize 9 & \scriptsize 10 & \textbf{\scriptsize Moyenne} & \textbf{\scriptsize Biais (µm)}\\\hline
5 mm &  &  &  &  &  &  &  &  &  &  &  & \\\hline
10 mm &  &  &  &  &  &  &  &  &  &  &  & \\\hline
15 mm &  &  &  &  &  &  &  &  &  &  &  & \\\hline
20 mm &  &  &  &  &  &  &  &  &  &  &  & \\\hline
25 mm &  &  &  &  &  &  &  &  &  &  &  & \\\hline
\end{tabularx}}\par{\small Biais = moyenne $-$ valeur de la cale. 1 µm = 0,001 mm.}\par
\titre{J'observe, je décide}
\begin{enumerate}[label=\textcolor{cyan}{\textbf{Q\arabic*.}},leftmargin=22pt]
\item Le biais est-il le même pour les cinq cales ? Monte-t-il, descend-il avec la longueur ?\reponse[2]
\item Ellistat juge-t-il la pente significative ? Et la justesse ?\reponse[1]
\item Biais constant : que faire ? Biais qui change avec la longueur : que faire ?\reponse[2]
\end{enumerate}
\begin{tcolorbox}[retenir]Un décalage constant se corrige par un réglage. Une erreur qui change avec la longueur est un défaut de \textbf{linéarité} : on restreint la plage ou on change d'instrument.\end{tcolorbox}
\vfill{\footnotesize\color{uimm}\textbf{Ressources de la base :} Données de secours.}
\end{document}